% Part: history
% Chapter: biographies 
% Section: alan-turing 

\documentclass[../../../include/open-logic-section]{subfiles}

\begin{document}

\olfileid{his}{bio}{tur} 

\olsection{آلان تورنغ}

وُلد آلان تورنغ بمايدا فيل من لندن يوم 23 يونيو 1912، ويُعد أبا لعلم
الحاسوب النظري. وقد مال منذ حداثته إلى العلوم الفيزيائية والرياضيات، غير
أن مدارسه لم تكن تعنى بما مال إليه، إذ كان جُلّ همها الأدب والدراسات
الكلاسيكية؛ فسوء تقديره عندهم، وكثر توبيخ معلّميه له.

\olphoto{turing-alan}{Alan Turing}

ثم دخل كلية كينغز بجامعة كامبريدج ودرس فيها الرياضيات. وفي سنة 1936 وضع
ما يسمى اليوم آلة تورنغ، يريد بذلك أن يحد مفهوم الدالة القابلة للحساب حداً
دقيقاً، وأن يبرهن امتناع حسم مسألة القرار. وكان ألونزو تشرش قد سبقه إلى
النتيجة وبرهنها بحساب لامبدا الذي أنشأه؛ ومع ذلك نُشرت رسالة تورنغ ونُبّه
فيها إلى نتيجة تشرش. ثم دعاه تشرش إلى برينستون، فأقام بها من سنة 1936 إلى
سنة 1938، ونال الدكتوراه تحت إشرافه.

ولم يشغله المنطق عن ميله القديم إلى العلوم الفيزيائية؛ فانتفع البريطانيون
بمهارته العملية حين خدم في قسم تحليل الشفرات ببلتشلي بارك أيام الحرب
العالمية الثانية. وكان له المقام الأول في حل شفرة إنيغما التي اتخذتها
البحرية الألمانية لمراسلاتها. وبخبرته في الإحصاء والتعمية، وبإدخال الآلات
الإلكترونية، استطاع الفريق أن ينشئ لفك الشفرة آلة سموها «بومب». وكان
لآرائه أثر كذلك في إنشاء كولوسوس، أول حاسوب إلكتروني قابل للبرمجة في
العالم، واستُعمل هو أيضاً ببلتشلي بارك في كسر شفرة لورنتس الألمانية.

وكان تورنغ مثلياً، على أنه خطب سنة 1942 جوان كلارك، وكانت زميلته ببلتشلي
بارك؛ ثم نقض الخطبة وأفضى إليها بأمره. وكانت له في حياته علاقات عاطفية
عدة، مع أن القانون البريطاني يومئذ كان يجعل الممارسة الجنسية بين الرجال
جرماً يعاقب عليه. وفي سنة 1952 سرق منزله رجل صاحب لعشيقه آنذاك؛ فلما رفع
تورنغ الأمر إلى الشرطة أقر بعلاقته المثلية، ظناً منه أن الحكومة ماضية إلى
إباحة تلك الأفعال. ولم يكن الأمر كما ظن، فوُجهت إليه تهمة الفحش الجسيم،
وخيّر فاختار، بدلاً من السجن، علاجاً بالهرمونات يخفض الشهوة الجنسية. ثم
وُجد ميتاً يوم 8 يونيو 1954 بجرعة مفرطة من السيانيد، والأرجح أنه قتل نفسه.
وأصدرت الملكة إليزابيث الثانية سنة 2013 عفواً ملكياً عنه.

\begin{reading}
ولسيرة مستوفاة لآلان تورنغ يُنظر~\citet{Hodges2014}. وقد كانت حياته وآثاره
باعثة على مسرحية~\emph{كسر الشفرة}، التي أُعدت للتلفزيون سنة 1996 ومثّل
ديريك جاكوبي فيها دور تورنغ. واستُلهم من حياته ومن أيام عمله ببلتشلي بارك،
استلهاماً غير ملتزم بالتفاصيل، فيلم~\emph{لعبة التقليد}، وهو من تمثيل
بنديكت كامبرباتش وكيرا نايتلي، وقد رُشح لجائزة الأوسكار
\citep{Imitation2014}.

وفي~\citet{Radiolab2012} حلقات مسموعة عدة في حياة تورنغ وعمله. ويمكن على
الشبكة مشاهدة وثائقي بي بي سي هورايزن~\emph{الحياة الغريبة للدكتور
  تورنغ وموته}~\citep{Sykes1992}. ويعرض~\citep{Theelen2012} شريطاً قصيراً
لآلة تورنغ تعمل وقد صنعت من مكعبات ليغو، احتفاء بمئوية مولده سنة 2012.

وأصل مقالة تورنغ في آلات تورنغ ومسألة القرار هو \citet{Turing1937}.
\end{reading}

\end{document}
